\begin{table}[htbp]
\centering
\footnotesize
\caption{Alternative environmental-good definition: OECD list vs.\ APEC list}
\label{tab:apec}
\begin{threeparttable}
\begin{tabular}{lcc}
\toprule
 & OECD list & APEC list \\
 & (used in this paper) & (54 codes) \\
\midrule
EP depth (WB) $\times$ Green & -0.00457 & 0.00501 \\
 & (0.00696) & (0.01266) \\
 & [0.512] & [0.693] \\
\addlinespace
EP count (TREND) $\times$ Green & 0.00181 & 0.00319 \\
 & (0.00182) & (0.00208) \\
 & [0.320] & [0.126] \\
\bottomrule
\end{tabular}
\begin{tablenotes}[flushleft]\footnotesize
\item \emph{Notes:} Baseline variant, collapsed panel.
\item \textbf{What this tests.} There is no single definition of ``environmental good''. This paper uses the OECD list; here it is replaced with the APEC list, a much narrower set agreed politically in 2012. If the result depended on which products we call green, the two columns would diverge.
\item Standard error in parentheses, $p$-value in brackets. $^{*}$ $p<0.10$, $^{**}$ $p<0.05$, $^{***}$ $p<0.01$.
\end{tablenotes}
\end{threeparttable}
\end{table}
